% LNCS-provided:
% \RequirePackage{amsthm}
% \newtheorem{definition}{Definition}
% \newtheorem{theorem}{Theorem}
% \newtheorem{lemma}{Lemma}

\AtEndPreamble{%
\theoremstyle{acmdefinition}
\newtheorem{remark}[theorem]{Remark}
\newtheorem{convention}[theorem]{Convention}}


\makeatletter
\newcommand*\FreezeMathAlpha[1]{%
  \@tempcnta=#1\relax
  \ifnum\@tempcnta>"7000\relax
    \advance\@tempcnta-"7000\relax
    \mathchardef#1=\@tempcnta
  \fi}
\AtBeginDocument{%
  \FreezeMathAlpha\Gamma  \FreezeMathAlpha\Delta   \FreezeMathAlpha\Theta
  \FreezeMathAlpha\Lambda \FreezeMathAlpha\Xi      \FreezeMathAlpha\Pi
  \FreezeMathAlpha\Sigma  \FreezeMathAlpha\Upsilon \FreezeMathAlpha\Phi
  \FreezeMathAlpha\Psi    \FreezeMathAlpha\Omega}
\makeatother

% https://tex.stackexchange.com/a/40185
\makeatletter
\let\@realLabel\label
\makeatother
\RequirePackage{amsmath}
\RequirePackage{mathtools}

\RequirePackage{xparse}

\RequirePackage{xspace}

% https://tex.stackexchange.com/a/521495
\makeatletter
\@ifdefinable\MakeTokenFromCsname{%
  \long\def\MakeTokenFromCsname#1#{%
    \romannumeral0\@MakeTokenFromCsname{#1}%
  }%
}%
\newcommand\@MakeTokenFromCsname[2]{%
  \expandafter\@@@MakeTokenFromCsname\expandafter{\csname#2\endcsname}{ #1}%
}%
\newcommand\@@@MakeTokenFromCsname[2]{#2#1}%
\makeatother

\makeatletter
\def\DMO#1{%
  \MakeTokenFromCsname\DeclareMathOperator{@#1}{#1}%
  \MakeTokenFromCsname\NewDocumentCommand{#1}{o}{\MakeTokenFromCsname{@#1}\IfValueT{##1}{_{##1}}}%
}
\makeatother

% A little slack for the occasional line that would otherwise stick into
% the margin: applied only to lines TeX cannot break any other way.
\setlength{\emergencystretch}{1em}

\DeclareMathOperator{\fc}{subj}
\DMO{dom}
\DMO{flv}
\DMO{fv}
\DMO{fchans}

\makeatletter

\newcommand{\acronym}[2]{\emph{#1}~(#2)}

\def\hyph-{\mbox{-}\nobreak\hspace{0pt}}
\def\pre#1-{$#1$\mbox{-}\nobreak\hspace{0pt}}
% \nobreak avoids linebreaks
% the 0-width space allows the suffix to be correctly hyphenated

\newcommand{\picalc}{%
    \texorpdfstring{\pre{\pi}-calculus}{pi-calculus}}
\newcommand{\lcalc}{%
    \texorpdfstring{\pre{\lambda}-calculus}{lambda-calculus}}

\makeatother

\newcommand{\Nat}{\mathbb{N}}

\RequirePackage{myproofs}
\RequirePackage{mathpartir}
% Inference rules: mathpartir \infer* inside proofrule(s); see proofrules.sty
\RequirePackage{proofrules}

% `\sepr` for regularly spaced, `\sepr*` for more spacing.
\makeatletter
\def\sepr{\@ifstar\@seprstar\@seprnostar}
\def\@seprnostar{\mathbin{\mbox{\large{$\mid$}}}}
\def\@seprstar{~~\@seprnostar~}
\makeatother

\def\<{\langle}
\def\>{\rangle}
\def\disjoint{\mathrel{\#}}
\NewDocumentCommand\set{m}{{\{#1\}}}
\RequirePackage{stmaryrd}
\NewDocumentCommand\mset{m}{{\lbag#1\rbag}}
\NewDocumentCommand\msetA{m}{{\lbag\begin{array}[t]{l} #1\rbag \end{array}}}
\NewDocumentCommand\update{mmm}{{#1\oplus[#2 \mapsto #3]}}

\let\ltxvdash\vdash
\newcommand{\types}{\mathrel{\color{black}\ltxvdash}}
\newcommand{\proves}{\mathrel{\color{black}\vDash}}

\newcommand{\ichoice}{\mathbin{{\color{cTy}\oplus}}}
\newcommand{\echoice}{\mathrel{{\color{cTy}\&}}}

\PassOptionsToPackage{dvipsnames,rgb}{xcolor}
\RequirePackage{xcolor}

% Semantic markup: sorts, metavariables and marked operators.  Everything
% below marks leaves only; \fFunc/\fType/\fProc survive as aliases.
\RequirePackage{lrlight}
\colorlet{cProc}{cPr}
\colorlet{cType}{cTy}
\colorlet{cFunc}{cTm}
\def\fCtx#1{{\color{cTm}\mathcal{#1}}}

\newcommand{\var}{\mathit}

\def\hole{{\color{cTm}[\cdot]}}

\RequirePackage{upgreek}
% Term keywords carry the term colour; \tyOp is their type-level twin.
\def\tOp#1{\mathinner{{\color{cTm}\mathsf{#1}}}}
\def\tyOp#1{\mathinner{{\color{cTy}\mathsf{#1}}}}
\def\tForce{\tOp{force}}
\def\tFork{\tOp{fork}}
\def\tRecv{\tOp{recv}}
\def\tSend{\tOp{send}}
\def\tWait{\tOp{wait}}
\def\tTerm{\tOp{term}}
\def\tClos{\tOp{clos}}
\def\tLink{\tOp{link}}
\def\tFixOp{\tOp{fix}}
\def\tForce{\tOp{force}}
\def\tInjOp{\tOp{inj}}
\def\tCaseOp{\tOp{case}}
\def\tIfOp{\tOp{if}}
\def\tUnit{{\color{cTm}\ast}}
\def\tyUnit{{\color{cTy}\ast}}
\let\tLVA\vr
\let\tLVB\vrB
\let\tIVX\ivx
\let\tIVY\ivy
\NewDocumentCommand\tAbs{mm}{{{\color{cTm}\uplambda}#1{\color{cTm}.}#2}}
\NewDocumentCommand\tPair{mm}{{\mathopen{{\color{cTm}\<}}#1\mathpunct{{\color{cTm},}}#2\mathclose{{\color{cTm}\>}}}}
\NewDocumentCommand\tBox{m}{{{\color{cTy}\Box}{#1}}}
\NewDocumentCommand\tApp{mm}{{#1\,#2}}
\NewDocumentCommand\tLet{soomm}{{\tOp{let}\IfBooleanTF{#1}{\tUnit}{\tPair{#2}{#3}}=#4\tOp{in}#5}}
\NewDocumentCommand\tFix{mmm}{{\tFixOp\,\tAbs{#1}{\tAbs{#2}{#3}}}}
\NewDocumentCommand\tSub{mmm}{{#1\mathopen{{\color{cTm}[}}#2{\color{cTm}/}#3\mathclose{{\color{cTm}]}}}}
\NewDocumentCommand\tLetI{mmm}{{\tOp{let}{\IExp{#1}}=#2\tOp{in}#3}}
\NewDocumentCommand\tLetS{mmm}{{\tOp{let}{{#1}}=#2\tOp{in}#3}}
\NewDocumentCommand\tFixI{mm}{{\tFixOp\,{#1}.{#2}}}
\NewDocumentCommand\tInj{mm}{{\tInjOp_{#1} #2}}
\NewDocumentCommand\tCase{mmmmm}{{\tCaseOp\,{#1}( #2.#3) (#4.#5)}}
\NewDocumentCommand\tIf{mmm}{{\tIfOp\,{#1}( #2)(#3)}}
\NewDocumentCommand\tSeq{mm}{{#1;#2}}



% Labelled arrows.  \stackon centred the label over a \longrightarrow of
% fixed length, so a long label stuck out at both ends.  \xrightarrow
% grows with its label instead, but bottoms out well short of
% \longrightarrow, so we pad the label to the width that makes the two
% agree: \xrightarrow's width is its label's width plus a fixed overhead,
% which we measure once.
\makeatletter
\newlength\lr@arw@min
\newlength\lr@arw@tmp
\AtBeginDocument{%
  \settowidth\lr@arw@min{$\longrightarrow$}%
  \settowidth\lr@arw@tmp{$\xrightarrow{\mathmakebox[\lr@arw@min]{}}$}%
  \addtolength\lr@arw@tmp{-\lr@arw@min}%  the overhead
  \addtolength\lr@arw@min{-\lr@arw@tmp}%  the label width that matches
  \ifdim\lr@arw@min<\z@ \lr@arw@min\z@ \fi}
\newcommand\lr@arw@lbl[1]{%
  \settowidth\lr@arw@tmp{$\scriptstyle#1$}%
  \ifdim\lr@arw@tmp<\lr@arw@min \lr@arw@tmp\lr@arw@min \fi%
  \mathmakebox[\lr@arw@tmp]{#1}}
\NewDocumentCommand\trans{mo}{%
\ifblank{#1}{\xrightarrow{}}{
  \xrightarrow{\lr@arw@lbl{#1}}
}%
  \IfValueT{#2}{^{#2}}}
% \redu carries no label, but still goes through \xrightarrow so that it
% is the same arrow, and the same length, as \trans with a short one.
\NewDocumentCommand\redu{o}{%
  \longrightarrow\IfValueT{#1}{^{#1}}}
\makeatother
% Step of the label-free reduction semantics (appendix on harmony).
\newcommand{\sred}{\mathrel{\rightsquigarrow}}
\def\irred{\operatorname{irred}}

% Label constructors own their own symbol only; the channels and values
% they mention arrive already marked.
\NewDocumentCommand\lRecv{mm}{{{#1}{?}{#2}}}
\NewDocumentCommand\lSend{mm}{{{#1}{!}{#2}}}
\NewDocumentCommand\lWait{m}{{{#1}{?}}}
\NewDocumentCommand\lTerm{m}{{{#1}{!}}}
\let\lClos\lTerm
\NewDocumentCommand\lFork{mm}{{(\nu{#1}){#2}}}
\NewDocumentCommand\lLink{mm}{{ {\leftrightarrow}\set{{#1}, {#2}}}}
\NewDocumentCommand\lLinkSingle{m}{{ {\leftrightarrow}\set{{#1}}}}

\RequirePackage{amssymb}
\NewDocumentCommand\lFunc{mm}{{{#1}\mathrel{{\color{cTy}\multimap}}{#2}}}
\NewDocumentCommand\lPair{mm}{{{#1}\mathrel{{\color{cTy}\otimes}}{#2}}}
\NewDocumentCommand\lSum{mm}{{{#1}\mathrel{{\color{cTy}+}}{#2}}}
\NewDocumentCommand\sRecv{mm}{{{\color{cTy}?}{#1}{\color{cTy}.}{#2}}}
\NewDocumentCommand\sSend{mm}{{{\color{cTy}!}{#1}{\color{cTy}.}{#2}}}
% \def\sWait{\tOp{Wait}}
% \def\sTerm{\tOp{Term}}
\def\sWait{\tyOp{Close}}
\def\sTerm{\tyOp{Close}}
\def\tEmpty{{\color{cTy}{}\cdot{}}}
\NewDocumentCommand\dual{m}{{\color{cTy}\overline{#1}}}
\def\ptype{\operatorname{ptype}}
\def\primitive{\operatorname{prim}}

\RequirePackage{stmaryrd}
\def\tRepl{\tOp{repl}}
\def\tReq{\tOp{req}}
\makeatletter
% \let\@replSymb\shortuparrow
\let\@replSymb!
\let\@reqSymb\shortdownarrow
\NewDocumentCommand\lRepl{mm}{{{\@replSymb}\lFork{#1}{#2}}}
\NewDocumentCommand\lReq{mm}{{{\@reqSymb}\lFork{#1}{#2}}}
\NewDocumentCommand\lUnres{m}{{{\@replSymb}#1}}
\NewDocumentCommand\IExp{m}{{{\@replSymb}#1}}
\makeatother

\RequirePackage{stmaryrd}
\makeatletter
\NewDocumentCommand\@lr{mmm}{{#1}_{#2}\llbracket{#3}\rrbracket}
\NewDocumentCommand\lrVal{mm}{\@lr{V}{#1}{#2}}
\NewDocumentCommand\lrTerm{mm}{\@lr{E}{#1}{#2}}
\NewDocumentCommand\lrProc{mm}{\@lr{P}{#1}{#2}}
\NewDocumentCommand\lrGSub{mm}{\@lr{G}{#1}{#2}}
% Script T, so the judgment letter does not collide with the
% metavariable T for types.
\NewDocumentCommand\lrTopo{mm}{\@lr{\mathcal{T}}{#1}{#2}}
\makeatother

% Σ ⊩_n (transition) : T  --  "this transition discharges its
% obligation at Σ ⊢ T, with continuations at index n".
\NewDocumentCommand\obl{mmm}{{#1} \mathrel{\Vdash_{#2}} {#3}}

\newcommand{\cupdot}{\mathbin{\mathaccent\cdot\cup}}
% The n-ary form of \cupdot is stmaryrd's \bigcupdot.

\newcommand{\ie}{i.e.,~} % comma
\newcommand{\eg}{e.g.~}  % no comma, according to the Handbook for Scholars

\newcommand{\st}{.\,}

\DeclareMathOperator{\unf}{unfold}
\newcommand{\sflipop}{\mathit{flip}}
\newcommand{\isdual}{\mathrel{\perp}}

\DeclareMathOperator{\sent}{sentc}
% \newcommand{\sent}{sentc}


%%% Local Variables:
%%% mode: latex
%%% TeX-master: "main"
%%% End:
